\documentclass[../DoAn.tex]{subfiles}
\begin{document}

Chương này giới thiệu tổng quan về đồ án tốt nghiệp với đề tài phát triển trò chơi hành động bắn súng Paradox Caliber. Nội dung chương trình bày lần lượt các vấn đề sau: (i) bối cảnh và sự cấp thiết của đề tài trong phần \ref{section:1.1}, (ii) mục tiêu cụ thể của đồ án trong phần \ref{section:1.2}, (iii) phạm vi và đối tượng người chơi hướng đến trong phần \ref{section:1.3}, (iv) định hướng giải pháp trong phần \ref{section:1.4}, và (v) bố cục tổng thể của báo cáo trong phần \ref{section:1.5}.

\section{Đặt vấn đề}
\label{section:1.1}

Ngành công nghiệp trò chơi điện tử trên toàn cầu đang trải qua giai đoạn tăng trưởng liên tục. Theo báo cáo thường niên của Newzoo \cite{Newzoo2024}, doanh thu thị trường game thế giới năm 2024 đạt khoảng 187,7 tỷ USD và dự kiến tiếp tục tăng trong các năm tới. Trong đó, phân khúc game trên nền tảng PC chiếm tỷ trọng đáng kể với khoảng 37,4 tỷ USD \cite{Newzoo2024}, phản ánh nhu cầu lớn của người dùng đối với các sản phẩm giải trí tương tác chất lượng cao trên nền tảng này. Trước đó, vào năm 2020, ngành game toàn cầu đã đạt mức doanh thu kỷ lục 179 tỷ USD, tăng 20\% so với năm 2019, cho thấy sức bền và khả năng tăng trưởng vượt trội so với nhiều ngành giải trí khác \cite{Witkowski2020}.

Trong số các thể loại trò chơi phổ biến, thể loại bắn súng hành động (Action Shooter) luôn duy trì vị trí hàng đầu về lượng người chơi và doanh thu. Các tựa game tiêu biểu như Control (Metascore 85/100 trên PC \cite{MetacriticControl}), Hitman: World of Assassination (Metascore 87/100 trên PC \cite{MetacriticHitman3}), và Max Payne (Metascore 89/100 trên PC \cite{MetacriticMaxPayne}) đã chứng minh sức hấp dẫn bền vững của thể loại bắn súng góc nhìn thứ ba (Third-Person Shooter -- TPS) qua nhiều thế hệ người chơi. Đặc biệt, những tựa game tích hợp các cơ chế sáng tạo như làm chậm thời gian (Bullet Time) trong Max Payne, hay sử dụng năng lực siêu nhiên để thao túng vật thể và không gian trong Control, đã nhận được phản hồi tích cực từ cả cộng đồng người chơi lẫn giới phê bình trên các trang đánh giá uy tín như Metacritic, bởi chúng mang lại trải nghiệm chiến thuật phong phú hơn so với các trò chơi bắn súng truyền thống. Tuy nhiên, trên thị trường hiện tại, số lượng sản phẩm kết hợp đồng thời các cơ chế bắn súng ẩn nấp, dịch chuyển không gian và thao túng thời gian trong một trò chơi TPS vẫn còn rất hạn chế.

Song song với nhu cầu của người chơi, sự phát triển nhanh chóng của thị trường game tại khu vực Đông Nam Á nói chung và Việt Nam nói riêng cũng mở ra triển vọng cho các nhà phát triển nội địa. Theo báo cáo của Niko Partners \cite{NikoPartners2024}, thị trường game Đông Nam Á năm 2024 có khoảng 270 triệu người chơi, trong đó Việt Nam đứng trong nhóm dẫn đầu về lượng người chơi game trên PC. Mặc dù phần lớn sản phẩm game phổ biến tại Việt Nam đến từ các nhà phát triển nước ngoài, nhu cầu và năng lực phát triển game nội địa đang dần được cải thiện. Việc một sinh viên hoặc nhà phát triển độc lập trong nước có thể tạo ra một sản phẩm game 3D hoàn chỉnh sẽ góp phần khẳng định tiềm năng của nguồn nhân lực công nghệ Việt Nam trong lĩnh vực này.

Bên cạnh đó, sự phát triển của các công cụ phát triển game thế hệ mới cũng mở ra cơ hội lớn cho các nhà phát triển độc lập (indie developer). Unreal Engine 5 với các công nghệ đồ họa tiên tiến như Nanite (Virtualized Geometry) và Lumen (Global Illumination) cho phép xây dựng môi trường 3D chất lượng cao mà không yêu cầu đội ngũ phát triển quy mô lớn. Đồng thời, hệ sinh thái tài nguyên mở ngày càng phong phú từ các nền tảng như Fab cung cấp hàng nghìn mô hình 3D, vật liệu và hoạt ảnh miễn phí, giúp giảm đáng kể chi phí sản xuất. Điều này đặt ra một bài toán thực tiễn: liệu một cá nhân có thể tận dụng các công cụ và tài nguyên hiện đại để phát triển một trò chơi hành động hoàn chỉnh với chi phí tối thiểu, đồng thời vẫn đảm bảo chất lượng trải nghiệm cho người chơi hay không.

Ngoài ra, từ góc độ học thuật, việc xây dựng một trò chơi điện tử hoàn chỉnh đòi hỏi sinh viên phải vận dụng tổng hợp nhiều lĩnh vực kiến thức: từ lập trình hệ thống, thiết kế kiến trúc phần mềm, xử lý toán học trong không gian 3D, đến thiết kế trải nghiệm người dùng và trí tuệ nhân tạo. Đây là một bài toán tổng hợp có giá trị thực tiễn cao, phù hợp với mục tiêu đào tạo của chương trình Khoa học máy tính, đồng thời cung cấp kinh nghiệm thực hành trực tiếp với một game engine mạnh mẽ và phổ biến trong ngành công nghiệp như Unreal Engine 5.

\section{Mục tiêu đề tài}
\label{section:1.2}

Trên thị trường game hành động bắn súng hiện tại, các sản phẩm phổ biến thường tập trung vào một trong hai hướng chính. Hướng thứ nhất là các tựa game bắn súng góc nhìn thứ ba với cơ chế ẩn nấp (Cover Shooter) như Gears of War (Metascore 94/100 \cite{MetacriticGearsOfWar}) hay The Division, nơi người chơi dựa vào hệ thống vật cản để chiến đấu. Hướng thứ hai là các tựa game khai thác cơ chế thao túng thời gian như Max Payne (Bullet Time, Metascore 89/100 \cite{MetacriticMaxPayne}) hay Superhot (thời gian chỉ trôi khi người chơi di chuyển). Một số sản phẩm khác như Portal và Splitgate (Metascore 71/100 \cite{MetacriticSplitgate}) lại khai thác cơ chế dịch chuyển không gian thông qua cổng không gian (Portal). Tuy nhiên, mỗi sản phẩm thường chỉ khai thác một cơ chế đơn lẻ, dẫn đến trải nghiệm chiến thuật bị giới hạn trong một khuôn khổ nhất định.

Về phía nhu cầu người chơi, cộng đồng game ngày càng tìm kiếm những trải nghiệm mới lạ, nơi nhiều cơ chế gameplay được kết hợp hài hòa để tạo ra chiều sâu chiến thuật. Theo báo cáo \textit{Essential Facts} của Entertainment Software Association (ESA) năm 2024 \cite{ESA2024}, khoảng 65\% người chơi cho biết họ ưu tiên lựa chọn trò chơi có cơ chế chơi sáng tạo và khác biệt so với các sản phẩm cùng thể loại. Điều này cho thấy tiềm năng của một sản phẩm có khả năng tích hợp đa cơ chế trong cùng một trải nghiệm.

Từ thực trạng trên, đồ án này đặt ra mục tiêu phát triển trò chơi Paradox Caliber -- một trò chơi bắn súng hành động góc nhìn thứ ba, trong đó người chơi có thể đồng thời sử dụng ba cơ chế cốt lõi: (i) bắn súng ẩn nấp (Cover Shooting) để chiến đấu dựa vào vật cản trong môi trường, (ii) cổng không gian (Portal) để dịch chuyển tức thời nhằm tạo lợi thế chiến thuật hoặc đánh úp đối phương, và (iii) làm chậm thời gian (Bullet Time) để né tránh đạn và phản ứng trong các tình huống nguy hiểm. Trò chơi được thiết kế với 07 màn chơi tuyến tính có bối cảnh và độ khó tăng dần, kèm theo hệ thống trí tuệ nhân tạo (AI) điều khiển kẻ thù có khả năng tự tìm chỗ nấp và phối hợp tấn công.

Mục tiêu quan trọng khác của đồ án là chứng minh tính khả thi của việc phát triển một trò chơi 3D hoàn chỉnh với chi phí tối thiểu. Cụ thể, đồ án tận dụng tối đa các tài nguyên 3D miễn phí từ nền tảng Fab và các nguồn mở khác, kết hợp với khả năng tự thiết kế các thành phần có thể tự thực hiện, nhằm giảm thiểu chi phí sản xuất mà vẫn đạt chất lượng đồ họa chấp nhận được.

\section{Phạm vi và đối tượng người chơi}
\label{section:1.3}

Paradox Caliber hướng đến nhóm người chơi nam và nữ trong độ tuổi từ 16 đến 30, là những người thường chơi các tựa game hành động và bắn súng trên nền tảng PC. Cụ thể, nhóm đối tượng chính bao gồm người chơi quen thuộc với thể loại Third-Person Shooter (TPS) và First-Person Shooter (FPS), đặc biệt là những người yêu thích các tựa game như Control, Hitman: World of Assassination, Max Payne hay Gears of War. Nhóm người chơi này thường có kinh nghiệm với các cơ chế bắn súng ẩn nấp và sẽ nhanh chóng tiếp cận được lối chơi cốt lõi của Paradox Caliber.

Ngoài ra, trò chơi cũng thu hút nhóm người yêu thích đề tài khoa học viễn tưởng (Sci-Fi). Bối cảnh tương lai của Paradox Caliber xoay quanh cuộc chiến chống lại siêu trí tuệ nhân tạo, kết hợp các yếu tố dịch chuyển không gian và thao túng thời gian, vốn là những chủ đề phổ biến và có sức hút lớn trong văn hóa Sci-Fi. Điều này giúp sản phẩm tiếp cận được cả những người chơi không thường xuyên chơi game bắn súng nhưng bị thu hút bởi cốt truyện và bối cảnh thế giới quan. Thêm vào đó, nhóm người chơi yêu thích game có tính chiến thuật cao -- như người hâm mộ thể loại Real-Time Strategy (RTS) hoặc Tactical Shooter -- cũng có thể tìm thấy sự hấp dẫn trong việc phối hợp ba cơ chế Portal, Bullet Time và Cover Shooting để vượt qua các tình huống chiến đấu phức tạp.

Về phân loại độ tuổi, Paradox Caliber dự kiến nhận xếp hạng ESRB ở mức T (Teen), phù hợp cho người chơi từ 13 tuổi trở lên, với mô tả nội dung bao gồm Violence (Bạo lực). Trò chơi không sử dụng máu thực tế (realistic blood), không chứa nội dung tình dục, và ngôn ngữ trong game không bao gồm các từ ngữ thô tục. Các yếu tố bạo lực được thể hiện ở mức độ cách điệu, phù hợp với phong cách hành động khoa học viễn tưởng. Do đó, sản phẩm phù hợp với cả đối tượng người chơi nam và nữ ở nhiều quốc gia, bao gồm cả thị trường châu Á và phương Tây.

Về phạm vi sản phẩm, Paradox Caliber là trò chơi một người chơi (Single-player) trên nền tảng Windows PC. Sản phẩm bao gồm 07 màn chơi tuyến tính hoàn chỉnh, hệ thống giao diện người dùng (Main Menu, HUD, Pause Menu), và bản build thực thi định dạng .exe. Mục tiêu về hiệu năng là đạt tối thiểu 60 FPS ở cấu hình máy tính tầm trung.

\section{Định hướng giải pháp}
\label{section:1.4}

Về thể loại, Paradox Caliber thuộc thể loại Third-Person Shooter (TPS) -- bắn súng hành động góc nhìn thứ ba, kết hợp các yếu tố Cover Shooter (bắn súng ẩn nấp). Đây là thể loại đã được kiểm chứng về mặt thị trường và có nền tảng cơ chế phù hợp để tích hợp thêm các yếu tố sáng tạo.

Về công nghệ, đồ án sử dụng Unreal Engine 5 làm nền tảng phát triển chính. Lý do lựa chọn Unreal Engine 5 dựa trên ba yếu tố: (i) hệ thống đồ họa tiên tiến với Nanite và Lumen cho phép tạo môi trường 3D chất lượng cao, (ii) hệ thống lập trình trực quan Blueprint phù hợp cho việc phát triển nhanh các nguyên mẫu (prototype) và triển khai logic game, và (iii) hệ sinh thái tài nguyên miễn phí phong phú từ Fab giúp giảm chi phí sản xuất. Bên cạnh đó, hệ thống AI Behavior Tree tích hợp sẵn trong Unreal Engine 5 cung cấp nền tảng để xây dựng trí tuệ nhân tạo cho kẻ thù mà không cần phát triển từ đầu.

Giải pháp cụ thể của đồ án là xây dựng một trò chơi hoàn chỉnh với ba cơ chế cốt lõi hoạt động đồng thời. Cơ chế thứ nhất là hệ thống bắn súng ẩn nấp, cho phép nhân vật dựa vào vật cản trong môi trường để chiến đấu. Cơ chế thứ hai là hệ thống cổng không gian (Portal), sử dụng các phép tính toán Vector và Rotator để thực hiện dịch chuyển tức thời mà không làm sai lệch hướng nhìn của người chơi. Cơ chế thứ ba là hệ thống làm chậm thời gian (Bullet Time), dựa trên thuật toán Global/Custom Time Dilation để tạo hiệu ứng chậm toàn cục hoặc cục bộ trong trò chơi. Ngoài ra, hệ thống AI kẻ thù được thiết kế dựa trên Behavior Tree với các trạng thái tuần tra, nhận diện và tấn công.

Đóng góp chính của đồ án bao gồm: (i) một sản phẩm trò chơi hoàn chỉnh với 07 màn chơi, tích hợp đồng thời ba cơ chế gameplay chưa phổ biến trên thị trường, (ii) giải pháp kỹ thuật cho bài toán dịch chuyển tức thời qua cổng không gian trong môi trường 3D mà không gây gián đoạn trải nghiệm, và (iii) minh chứng cho tính khả thi của việc phát triển trò chơi 3D chất lượng bởi một cá nhân với chi phí tối thiểu nhờ tận dụng công cụ và tài nguyên mở.

\section{Bố cục đồ án}
\label{section:1.5}

Phần còn lại của báo cáo đồ án tốt nghiệp này được tổ chức như sau.

Chương 2 trình bày quá trình khảo sát và phân tích yêu cầu cho trò chơi Paradox Caliber. Nội dung chương bao gồm việc khảo sát các sản phẩm game hành động bắn súng tương tự trên thị trường, phân tích và so sánh ưu nhược điểm của từng sản phẩm, từ đó xác định các tính năng cốt lõi cần phát triển và đưa ra mục đích cũng như định hướng sử dụng cho trò chơi.

Trong Chương 3, em trình bày các công nghệ và nền tảng lý thuyết nổi bật được sử dụng trong đồ án. Chương này tập trung vào các nguyên lý toán học trong không gian 3D phục vụ cho cơ chế Portal, thuật toán Time Dilation cho cơ chế Bullet Time, và nền tảng lý thuyết về Behavior Tree cho hệ thống trí tuệ nhân tạo điều khiển kẻ thù.

Chương 4 trình bày chi tiết về thiết kế trò chơi Paradox Caliber, bao gồm tổng quan trò chơi, lối chơi, các cơ chế cốt lõi, hệ thống điều khiển, cốt truyện, thiết kế kiến trúc phần mềm, thiết kế chi tiết các lớp và hệ thống, cũng như quá trình xây dựng 07 màn chơi với bối cảnh từ sa mạc, khu công nghiệp, đô thị đến không gian siêu thực.

Chương 5 trình bày kết quả thực nghiệm và đánh giá sản phẩm. Nội dung chương bao gồm các thông số kỹ thuật đạt được, kết quả kiểm thử chức năng trên từng màn chơi, đánh giá hiệu năng về mặt khung hình (FPS), và quy trình đóng gói sản phẩm thành bản thực thi trên nền tảng Windows.

Chương 6 tập trung vào các giải pháp và đóng góp nổi bật mà em đã thực hiện trong suốt quá trình phát triển đồ án. Các giải pháp được phân tích bao gồm thuật toán xử lý dịch chuyển qua cổng không gian mà không gây sai lệch hướng nhìn, giải pháp tối ưu hiệu năng đồ họa, và cách thức tích hợp đa cơ chế gameplay trong cùng một hệ thống.

Chương 7 đưa ra kết luận tổng thể về đồ án và trình bày các hướng phát triển trong tương lai, bao gồm khả năng mở rộng thêm màn chơi, bổ sung chế độ nhiều người chơi, và cải thiện hệ thống trí tuệ nhân tạo cho kẻ thù.

Chương này đã giới thiệu bối cảnh thực tiễn của ngành công nghiệp game, nhu cầu thị trường đối với các trò chơi bắn súng hành động có cơ chế sáng tạo, cũng như mục tiêu và phạm vi cụ thể của đồ án Paradox Caliber. Trên cơ sở các vấn đề và mục tiêu đã xác định, Chương 2 tiếp theo sẽ trình bày chi tiết quá trình khảo sát hiện trạng các sản phẩm tương tự và phân tích yêu cầu cho trò chơi.

\end{document}
